\scanpage{156}
qu'il est \uline{surjectif} et \uline{radiciel}, donc que pour tt $y\in Y$, la fibre géométrique de $P/H$ \deleted{\unsure{s'envoie}} en $y$ s'envoie \unsure{bijectivemt|bijectivement} sur celle de $X$ en $y$. Or c'est encore un pt trivial, d'après l'explication précédente.

% slanted
\margin{pts \unsure{géométriques} de \unsure{base|une} de $P/H$ $=$ \deleted{\unsure{pts géo}} \unsure{classes|orbites} sous $G$ dans les pts géom. de $P$.}

Évidemt (iii bis) $\Rightarrow$ (iii), et (iii) $\Rightarrow$ (i) d'après un \uline{lemme} précédent, déjà utilisé. D'où l'équivalence de (i) (iii bis) (iii).

De plus (iii\deleted{\unsure{bis}}) $\Rightarrow$ (ii bis), comme on voit en prenant $Y_1=P$, et utilisant la formule $(X\times\deleted{\unsure{H}}G)/H\simeq X\times G/H$, et (ii bis) $\Rightarrow$ (ii). Enfin, (ii \deleted{bis}) $\Rightarrow$ (i) est immédiat \deleted{[car la condition \unsure{sur les}} car si $X$ est défini par un faisceau

\margin{\uline{Cor}.}

coh. $\underline A$ sur $Y$, $X_1$ le sera par son image inverse $\underline A_1$. Si $\underline A_1$ est loc. libre, il en est de $=$ de $\underline A$ comme bien connu. Si les fibres de $X_1$ sur $Y_1$ sont séparables \inserted{(déduites de celles de $X$ sur $Y$ par ext. des corps de base)}, il en est de $=$ de celles de $X$ sur $Y$, \deleted{\unsure{Donc}} Donc si $X_1$ est étale sur $Y_1$, $X$ l'est sur $Y$.

\uline{Cor} Si $X$ est connexe, alors dans (iii) on peut prendre $P$ \uline{connexe}.

% archive sheet number: 155

